% OmniLaTeX Advanced Beamer Showcase
% Compile with: latexmk -lualatex main.tex
%
% Demonstrates: overlays & animations, blocks & environments, columns &
% layouts, TikZ integration, navigation & structure, advanced formatting,
% code listings, math typesetting, footnotes & references, appendix slides.

\documentclass[
  doctype=beamer,
  beameraspectratio=169,
  beamertheme=Madrid,
  institution=generic,
]{omnilatex}

\usecolortheme{whale}

\usepackage{tikz}
\usetikzlibrary{arrows.meta,positioning,calc,decorations.pathreplacing}
\usepackage{listings}
\usepackage{xcolor}
\usepackage{amsmath,amssymb}

\lstset{
  basicstyle=\ttfamily\small,
  keywordstyle=\color{blue}\bfseries,
  commentstyle=\color{green!50!black}\itshape,
  stringstyle=\color{red!70!black},
  numbers=left,
  numberstyle=\tiny\color{gray},
  stepnumber=1,
  frame=single,
  breaklines=true,
  showstringspaces=false,
  tabsize=2,
  language=Python,
  morekeywords={self,True,False,None},
}

\newcommand{\checkmark}{\text{\(\checkmark\)}}

% ── Presentation metadata ──────────────────────────
\title[Advanced Beamer Features]{Advanced Beamer Features\\
with OmniLaTeX and LuaLaTeX}
\subtitle{A Comprehensive Showcase}
\author[Wyatt Au]{Wyatt Au}
\institute[OmniLaTeX Project]{OmniLaTeX Project\\\texttt{https://github.com/WyattAu/OmniLaTeX-template}}
\date{June 2026}

% ── Bibliography ────────────────────────────────────
\usepackage[style=ieee]{biblatex}
\addbibresource{../../bib/bibliography.bib}

\begin{document}

%% ════════════════════════════════════════════════════
%% FRAME: Title Page
%% ════════════════════════════════════════════════════

\begin{frame}[plain]
  \titlepage
\end{frame}

%% ════════════════════════════════════════════════════
%% FRAME: Table of Contents
%% ════════════════════════════════════════════════════

\begin{frame}{Outline}
  \tableofcontents[hideallsubsections]
\end{frame}

%% ────────────────────────────────────────────────────
\section{Overlays and Animations}
%% ────────────────────────────────────────────────────

\begin{frame}{Conditional Content}
  \begin{block}{\texttt{\textbackslash only} vs.\ \texttt{\textbackslash uncover} vs.\ \texttt{\textbackslash visible}}
    \begin{itemize}
      \item \texttt{\textbackslash only<2->{}}: Content appears from slide 2 onward; space reserved only when visible.
      \item \texttt{\textbackslash uncover<2-3>{}}: Content appears on slides 2--3; space always reserved.
      \item \texttt{\textbackslash visible<2->{}}: Like \texttt{\textbackslash uncover}, but invisible content is also hidden from PDF.
    \end{itemize}
  \end{block}

  \only<2->{\begin{alertblock}{This block appears on slide 2 onward}
    Demonstrates \texttt{\textbackslash only<2->{}}.
  \end{alertblock}}

  \uncover<3-4>{\begin{exampleblock}{This block is visible on slides 3--4}
    Demonstrates \texttt{\textbackslash uncover<3-4>{}}.
  \end{exampleblock}}
\end{frame}

\begin{frame}{Step-by-Step Reveals with \texttt{\textbackslash pause}}
  Step 1: First piece of content.\pause
  Step 2: Second piece of content.\pause
  Step 3: Third piece of content.\pause
  Step 4: Final piece of content.

  \vspace{1em}
  \begin{block}{Note}
    \texttt{\textbackslash pause} inserts a slide break at each invocation.
  \end{block}
\end{frame}

\begin{frame}{Switching and Temporal Effects}
  \begin{columns}[T]
    \column{0.5\textwidth}
    \begin{block}{\texttt{\textbackslash alt}}
      \alt<1|2>{First option}{Second option}
    \end{block}

    \column{0.5\textwidth}
    \begin{block}{\texttt{\textbackslash temporal}}
      \temporal<1|2|3>{Before}{At current}{After}
    \end{block}
  \end{columns}

  \vspace{1em}
  \begin{alertblock}{Try navigating!}
    Use \texttt{<} and \texttt{>} keys or click to step through overlays.
  \end{alertblock}
\end{frame}

\begin{frame}{Incremental Item Reveals}
  \begin{itemize}
    \item<1-> First item: Always visible.
    \item<2-> Second item: Appears on slide 2.
    \item<3-> Third item: Appears on slide 3.
    \item<4-> Fourth item: Appears on slide 4.
    \item<5-> Fifth item: Appears on slide 5.
  \end{itemize}

  \vspace{0.5em}
  \begin{exampleblock}{Custom bullets}
    \begin{itemize}
      \item[\checkmark] Completed task
      \item[\checkmark] Another done task
      \item[---] Pending task
    \end{itemize}
  \end{exampleblock}
\end{frame}

%% ────────────────────────────────────────────────────
\section{Blocks and Environments}
%% ────────────────────────────────────────────────────

\begin{frame}{Block Environments}
  \begin{columns}[T]
    \column{0.48\textwidth}
    \begin{block}{Standard Block}
      General-purpose highlighted content.
    \end{block}
    \begin{alertblock}{Alert Block}
      Warnings and critical information.
    \end{alertblock}

    \column{0.48\textwidth}
    \begin{exampleblock}{Example Block}
      Demonstrations and examples.
    \end{exampleblock}
    \begin{block}{Custom Colored Block}
      \setbeamercolor{block body}{bg=blue!10}
      Adapt block colors per-block.
    \end{block}
  \end{columns}
\end{frame}

\begin{frame}{Theorem and Proof Environments}
  \begin{columns}[T]
    \column{0.5\textwidth}
    \begin{block}{Theorem 1 (Cauchy--Schwarz)}
      For vectors $u, v$ in an inner product space:
      \[
        |\langle u, v \rangle|^2 \leq \langle u, u \rangle \cdot \langle v, v \rangle
      \]
    \end{block}

    \column{0.5\textwidth}
    \begin{block}{Proof}
      Consider the non-negativity of $\|u - \lambda v\|^2 \geq 0$ for all $\lambda \in \mathbb{R}$.
      Expanding and minimizing over $\lambda$ yields the inequality.
    \end{block}
  \end{columns}

  \begin{block}{Lemma 2}
    If $f$ is continuous on $[a,b]$ and differentiable on $(a,b)$, then there exists $c \in (a,b)$ such that $f'(c) = \frac{f(b)-f(a)}{b-a}$.
  \end{block}
\end{frame}

\begin{frame}[fragile]{Code Listings in Beamer}
  \begin{block}{Python: Fibonacci with Memoization}
    \begin{lstlisting}[language=Python]
from functools import lru_cache

@lru_cache(maxsize=None)
def fib(n):
    if n < 2:
        return n
    return fib(n - 1) + fib(n - 2)

# Compute first 10 Fibonacci numbers
print([fib(i) for i in range(10)])
    \end{lstlisting}
  \end{block}
\end{frame}

%% ────────────────────────────────────────────────────
\section{Columns and Layouts}
%% ────────────────────────────────────────────────────

\begin{frame}{Side-by-Side Comparison}
  \begin{columns}[T]
    \column{0.48\textwidth}
    \begin{block}{Approach A: Procedural}
      \begin{enumerate}
        \item Read input file
        \item Parse line by line
        \item Transform tokens
        \item Write output file
      \end{enumerate}
    \end{block}

    \column{0.48\textwidth}
    \begin{block}{Approach B: Declarative}
      \begin{enumerate}
        \item Define grammar rules
        \item Specify transformations
        \item Apply pipeline
        \item Collect results
      \end{enumerate}
    \end{block}
  \end{columns}

  \vspace{1em}
  \rule{1\textwidth}{0.4pt}

  \begin{center}
    \textbf{\color{OmniPrimaryColor}Both approaches yield equivalent results.}
  \end{center}
\end{frame}

\begin{frame}{Multi-Column with Minipages}
  \begin{minipage}{0.3\textwidth}
    \begin{block}{Column 1}
      First column content with some text to demonstrate the layout.
    \end{block}
  \end{minipage}\hfill
  \begin{minipage}{0.3\textwidth}
    \begin{block}{Column 2}
      Second column content. Minipages allow fine-grained width control.
    \end{block}
  \end{minipage}\hfill
  \begin{minipage}{0.3\textwidth}
    \begin{block}{Column 3}
      Third column content. Three equal-width columns with \texttt{\textbackslash hfill} spacing.
    \end{block}
  \end{minipage}
\end{frame}

%% ────────────────────────────────────────────────────
\section{TikZ Integration}
%% ────────────────────────────────────────────────────

\begin{frame}{Animated TikZ Diagram}
  \begin{center}
    \begin{tikzpicture}[
      node distance=1.5cm,
      block/.style={rectangle, draw, fill=blue!5, minimum width=2cm,
                    minimum height=0.8cm, align=center, rounded corners},
      highlight/.style={rectangle, draw, fill=red!20, minimum width=2cm,
                        minimum height=0.8cm, align=center, rounded corners, thick},
    ]
      % Step 1: Input node
      \node[block] (input) {Input};
      % Step 2: Process node (hidden initially)
      \only<2->{\node[block, right=of input] (process) {Process};}
      % Step 3: Decision node (hidden initially)
      \only<3->{\node[block, right=of process] (decision) {Decision};}
      % Step 4: Output node (hidden initially)
      \only<4->{\node[block, right=of decision] (output) {Output};}

      % Arrows
      \only<2->{\draw[->,>=stealth,thick] (input) -- (process);}
      \only<3->{\draw[->,>=stealth,thick] (process) -- (decision);}
      \only<4->{\draw[->,>=stealth,thick] (decision) -- (output);}

      % Highlight on specific slides
      \only<3>{\node[highlight, right=of process] (decision) {Decision};}
    \end{tikzpicture}
  \end{center}

  \vspace{0.5em}
  \begin{center}
    \small Step-by-step construction of a flow diagram using \texttt{\textbackslash only}
  \end{center}
\end{frame}

\begin{frame}{TikZ with Highlighted Nodes}
  \begin{center}
    \begin{tikzpicture}[
      every node/.style={circle, draw, minimum size=1cm},
      highlight/.style={fill=red!30, draw=red!70!black, thick},
    ]
      \node (a) at (0,0) {$A$};
      \node (b) at (2,0) {$B$};
      \node (c) at (1,1.73) {$C$};
      \node (d) at (1,-1.73) {$D$};

      \draw[->,>=stealth] (a) -- (b);
      \draw[->,>=stealth] (b) -- (c);
      \draw[->,>=stealth] (c) -- (a);
      \draw[->,>=stealth] (a) -- (d);
      \draw[->,>=stealth] (d) -- (b);

      \only<2>{\node[highlight] at (0,0) {$A$};}
      \only<3>{\node[highlight] at (2,0) {$B$};}
      \only<4>{\node[highlight] at (1,1.73) {$C$};}
      \only<5>{\node[highlight] at (1,-1.73) {$D$};}
    \end{tikzpicture}
  \end{center}

  \begin{center}
    \small Highlight individual nodes with \texttt{\textbackslash only<slide>\{\textbackslash node[highlight]...\};}
  \end{center}
\end{frame}

%% ────────────────────────────────────────────────────
\section{Navigation and Structure}
%% ────────────────────────────────────────────────────

\begin{frame}{Table of Contents}
  \begin{columns}[T]
    \column{0.5\textwidth}
    \begin{block}{Full Table of Contents}
      \tableofcontents[hideallsubsections]
    \end{block}

    \column{0.5\textwidth}
    \begin{block}{Current Section Only}
      \tableofcontents[currentsection, hideallsubsections]
    \end{block}
  \end{columns}
\end{frame}

\subsection{Subsection Example}
\begin{frame}{Subsection Content}
  This frame belongs to a subsection, demonstrating hierarchical structure.
  \begin{block}{Key Point}
    Use \texttt{\textbackslash section} and \texttt{\textbackslash subsection} to organize your presentation.
  \end{block}
\end{frame}

%% ────────────────────────────────────────────────────
\section{Advanced Formatting}
%% ────────────────────────────────────────────────────

\begin{frame}{Emphasis and Alignment}
  \begin{columns}[T]
    \column{0.5\textwidth}
    \begin{block}{Text Emphasis}
      \begin{itemize}
        \item \textbf{Bold text} for strong emphasis
        \item \textit{Italic text} for terms
        \item \texttt{Monospace} for code references
        \item \textcolor{red}{Colored text} for warnings
        \item \textbf{\color{blue!70!black}Combined} emphasis
      \end{itemize}
    \end{block}

    \column{0.5\textwidth}
    \begin{block}{Custom Bullets}
      \begin{itemize}
        \item[\checkmark] Completed task
        \item[$\times$] Failed task
        \item[$\diamond$] Optional task
        \item[$\circ$] Pending task
        \item[$\bullet$] Standard bullet
      \end{itemize}
    \end{block}
  \end{columns}

  \vspace{1em}
  \begin{center}
    \rule{1\textwidth}{0.4pt}
  \end{center}
  \begin{center}
    \raggedright
    This text is \texttt{\textbackslash raggedright} aligned, stopping at the margin.
  \end{center}
\end{frame}

%% ────────────────────────────────────────────────────
\section{Code Listings}
%% ────────────────────────────────────────────────────

\begin{frame}[fragile]{Python with Line Highlighting}
  \begin{block}{Highlighted Lines}
    \begin{lstlisting}[language=Python, highlightlines={3-5}]
import numpy as np

def normalize(x):
    """Normalize array to [0, 1] range."""
    x_min, x_max = x.min(), x.max()
    return (x - x_min) / (x_max - x_min)

data = np.random.rand(100)
result = normalize(data)
    \end{lstlisting}
  \end{block}
\end{frame}

\begin{frame}[fragile]{LaTeX Code Listing}
  \begin{block}{Document Class Usage}
    \begin{lstlisting}[language={[LaTeX]TeX}, basicstyle=\ttfamily\scriptsize]
\documentclass[
  doctype=thesis,
  institution=tuhh,
  language=german,
  font=libertinus,
]{omnilatex}
    \end{lstlisting}
  \end{block}
\end{frame}

%% ────────────────────────────────────────────────────
\section{Math in Beamer}
%% ────────────────────────────────────────────────────

\begin{frame}{Step-by-Step Equation Building}
  \onslide<2->{%
    \begin{equation}
      \int_{-\infty}^{\infty} e^{-x^2} \, dx = \sqrt{\pi}
    \end{equation}
  }

  \vspace{1em}
  \onslide<3->{%
    \begin{block}{Gaussian Integral}
      \[
        I = \int_{-\infty}^{\infty} e^{-x^2} \, dx
        \implies I^2 = \pi
        \implies I = \sqrt{\pi}
      \]
    \end{block}
  }
\end{frame}

\begin{frame}{Matrix with Highlighted Elements}
  \[
    A = \begin{pmatrix}
      \only<2>{\color{red}\mathbf{a_{11}}} \only<3->{a_{11}} &
      \only<3>{\color{blue}\mathbf{a_{12}}} \only<4->{a_{12}} &
      a_{13} \\
      a_{21} &
      \only<4>{\color{green!50!black}\mathbf{a_{22}}} \only<5->{a_{22}} &
      a_{23} \\
      a_{31} &
      a_{32} &
      \only<5>{\color{purple}\mathbf{a_{33}}} \only<6->{a_{33}}
    \end{pmatrix}
  \]

  \vspace{1em}
  \onslide<6->
  \begin{block}{Determinant Expansion}
    \[
      \det(A) = a_{11}(a_{22}a_{33} - a_{23}a_{32})
              - a_{12}(a_{21}a_{33} - a_{23}a_{31})
              + a_{13}(a_{21}a_{32} - a_{22}a_{31})
    \]
  \end{block}
\end{frame}

%% ────────────────────────────────────────────────────
\section{Footnotes and References}
%% ────────────────────────────────────────────────────

\begin{frame}{Footnotes and Citations}
  \begin{block}{Timed Footnotes}
    This content is always visible.\footnote<2->{This footnote appears on slide 2 and onward.}

    Additional context.\footnote<3->{This footnote appears on slide 3 onward.}
  \end{block}

  \vspace{1em}
  \begin{block}{Citations}
    \begin{itemize}
      \item Knuth's foundational work\cite{knuthTeXbook1986} on typesetting.
      \item Einstein's relativity paper\cite{einstein1905}.
      \item McKinney's Python data analysis guide\cite{mckinneyPythonDataAnalysis2018}.
    \end{itemize}
  \end{block}
\end{frame}

%% ════════════════════════════════════════════════════
%% FRAME: Summary
%% ════════════════════════════════════════════════════

\begin{frame}{Summary}
  \begin{enumerate}
    \item Overlays and animations enable step-by-step content delivery.
    \item Block environments provide structured, colored content areas.
    \item Columns and minipages support flexible multi-column layouts.
    \item TikZ diagrams can be animated and highlighted per-slide.
    \item Section/subsection structure with dynamic table of contents.
    \item Code listings with line highlighting for source code display.
    \item Math environments support progressive equation building.
    \item Timed footnotes and bibliography citations enhance references.
  \end{enumerate}
\end{frame}

%% ════════════════════════════════════════════════════
%% FRAME: Thank You
%% ════════════════════════════════════════════════════

\begin{frame}[plain,c]
  \vspace{2cm}
  {\huge Thank You}\\[1cm]
  {\large Questions?}\\[2cm]
  {\footnotesize
  \texttt{https://github.com/WyattAu/OmniLaTeX-template}\\
  \texttt{omnilatex-template.wyattau.com}
  }
\end{frame}

%% ════════════════════════════════════════════════════
%% APPENDIX: Backup Slides
%% ════════════════════════════════════════════════════

\appendix

\begin{frame}[allowframebreaks]{Appendix: Full Doctype Reference}
  \begin{table}
    \centering
    \begin{tabular}{l l l}
      \hline
      Doctype         & Base Class & Aliases \\
      \hline
      thesis          & scrbook    & theses \\
      dissertation    & scrbook    & dissertations \\
      book            & scrbook    & --- \\
      dictionary      & scrbook    & dictionaries, lexicon, lexicons \\
      manual          & scrreprt   & manuals, guide, guides, handbook, handbooks \\
      technicalreport & scrreprt   & report, reports, technical-report, etc. \\
      standard        & scrreprt   & standards \\
      patent          & scrreprt   & patents \\
      article         & scrartcl   & articles, paper, papers \\
      cv              & scrartcl   & resume, resumes, curriculumvitae \\
      cover-letter    & scrartcl   & coverletter \\
      journal         & scrartcl   & journals, magazine, magazines \\
      poster          & scrartcl   & posters \\
      beamer          & beamer     & --- \\
      \hline
    \end{tabular}
  \end{table}
  (Full list: 27 canonical profiles with 55+ aliases)
\end{frame}

\begin{frame}{Appendix: Detailed Derivation}
  \begin{block}{Cauchy--Schwarz Inequality Proof Sketch}
    For $u, v \in V$ over field $\mathbb{F}$:
    \[
      t \mapsto \|u + tv\|^2 = \langle u + tv, u + tv \rangle \geq 0
    \]
    Expanding:
    \[
      \|u\|^2 + 2t \operatorname{Re}\langle u, v \rangle + t^2 \|v\|^2 \geq 0
    \]
    This quadratic in $t$ is non-negative, so its discriminant satisfies:
    \[
      \Delta = 4(\operatorname{Re}\langle u, v \rangle)^2 - 4\|u\|^2\|v\|^2 \leq 0
    \]
    Therefore $|\langle u, v \rangle|^2 \leq \|u\|^2 \|v\|^2$.
  \end{block}
\end{frame}

\begin{frame}{Appendix: Overlay Reference Table}
  \begin{table}
    \centering
    \small
    \begin{tabular}{l l}
      \hline
      Command & Behavior \\
      \hline
      \texttt{\textbackslash only<n>{}} & Show only on slide $n$; no space reserved \\
      \texttt{\textbackslash only<n->{}} & Show from slide $n$ onward \\
      \texttt{\textbackslash only<n-m>{}} & Show on slides $n$ through $m$ \\
      \texttt{\textbackslash uncover<n->{}} & Like \texttt{\textbackslash only}, space always reserved \\
      \texttt{\textbackslash visible<n->{}} & Like \texttt{\textbackslash uncover}, hidden content not in PDF \\
      \texttt{\textbackslash alt<n|n>{}} & Switch between two alternatives \\
      \texttt{\textbackslash temporal<n|n|n>{}} & Before, at, after overlays \\
      \texttt{\textbackslash pause} & Insert slide break \\
      \texttt{\textbackslash item<n->} & Incremental item reveal \\
      \texttt{\textbackslash footnote<n->{}} & Timed footnote \\
      \hline
    \end{tabular}
  \end{table}
\end{frame}

\end{document}
